% ============================================================
%  ch04a：§4 开头、4.1 基点与 Weil 群、4.2 伴随商（原文 pp.25–27）
%  起止：自 “4. Grothendieck–Springer theory for LS^arthm” 一节开头，
%        至 “4.3. Non-resonance” 标题之前（§4.3 起属 ch04b）。
% ============================================================
\section{$\LS^{\arthm}$ 的 Grothendieck--Springer 理论}

本节的目的是证明定理 3.5.2.1。由于该定理完全落在谱侧，本节通篇默认在域 $\e$ 上工作。

\subsection{基点与 Weil 群记号}

下面我们取定一个几何点 $x_0 \in X(k)$ 作为基点，它将在我们的基本群中充当基点的角色；这里提醒一句：$k = \Fb$。

我们建议读者对自己宽容一些，假定 $x_0$ 定义在 $\F_q$ 上；此时 §4.1 余下的材料基本上可以一概略去不看。

话虽如此，为使 $X_0$ 没有有理点的情形也能覆盖，我们在这里收录了一些技术性材料。

\pnum{}令 $\widetilde{X}$ 表示 $X$ 以 $x_0$ 为基点的万有覆盖；按定义，$\widetilde{X}$ 连通、在 $X$ 上 pro-有限 étale，配有 $x_0$ 的一个提升 $\widetilde{x}_0$，并在所有这类数据中是初始的。注意 $\widetilde{X}$ 同时也是 $X_0$ 的万有覆盖，因此存在 $\pi_1^{\et}(X_0)$（$= \pi_1^{\et}(X_0,x_0)$）在 $\widetilde{X}$ 上的一个同义反复的作用（它把 $\widetilde{X}$ 实现为 $X$ 上的 $\pi_1^{\et}(X)$-torsor、以及 $X_0$ 上的 $\pi_1^{\et}(X_0)$-torsor，其中每个 torsor 都理解为对 pro-étale 拓扑局部平凡）。

令 $\Frob_X : X \to X$ 表示几何 Frobenius 映射。我们一次性选定点 $\Frob_X(x_0)$ 到 $\widetilde{X}$ 的一个提升。容易看出，存在唯一的映射 $\widetilde{\Frob}_X$ 使下图交换
\[
\begin{tikzcd}[column sep=large, row sep=large]
  \widetilde{X} \arrow[r, "\widetilde{\Frob}_X"] \arrow[d, "\pi"'] & \widetilde{X} \arrow[d, "\pi"] \\
  X \arrow[r, "\Frob_X"] & X
\end{tikzcd}
\]
并把 $\widetilde{x}_0$ 送到我们选定的 $\Frob_X(x_0)$ 的提升（今后记作 $\widetilde{\Frob}_X(\widetilde{x}_0)$）。

这一选择还（与此相关地）定义了 $\widehat{\ZZ}$ 在 $\mathrm{Sch}_{/X_0}$ 中的 $\widetilde{X}$ 上的一个作用；生成元 $-1 \in \widehat{\ZZ}$ 的逆\footnote{我们指出，按本文的约定，当我们把 $\widehat{\ZZ} \simeq \mathrm{Gal}(\F_q)$ 等同起来时，生成元 $-1 \in \widehat{\ZZ}$ 的逆对应的是\emph{算术} Frobenius。}通过一个映射 $\gamma : \widetilde{X} \to \widetilde{X}$ 作用，该映射的特征是：在 $k$ 上是 Frobenius 半线性的，并且 $\gamma \circ \widetilde{\Frob}_X$ 是 $\widetilde{X}$ 的绝对 Frobenius。

这份数据定义了映射 $\mathcal{W}_X \to \ZZ$ 的一个分裂。我们令 $F \in \mathcal{W}_X$ 表示 $1 \in \ZZ$ 在该分裂下的像。按定义，$F^{-1} \in \mathcal{W}_X \subseteq \pi_1^{\et}(X_0)$ 通过上面记作 $\gamma$ 的映射作用在 $\widetilde{X}$ 上。

对 $g \in \pi_1^{\et}(X)$，我们令 ${}^F g := \Ad_F(g)$。注意点 $\widetilde{\Frob}_X(\widetilde{x}_0)$ 的选取给出了同构 $\pi_1^{\et}(X,\Frob_X(x_0)) \simeq \pi_1^{\et}(X,x_0)$，而复合
\[
  \pi_1^{\et}(X,x_0) \xrightarrow{\ \pi_1^{\et}(\Frob_X)\ } \pi_1^{\et}(X,\Frob_X(x_0)) \simeq \pi_1^{\et}(X,x_0)
\]
就是映射 $g \mapsto {}^F g$\footnote{确实，按 $\pi_1^{\et}(\Frob_X)(g)$ 的定义，有 $\widetilde{\Frob}_X(g\widetilde{x}_0) = \pi_1^{\et}(\Frob_X)(g) \cdot \widetilde{\Frob}_X(\widetilde{x}_0)$。
  又按 $F = \gamma^{-1}$ 的定义，有 $F^{-1}\cdot\widetilde{\Frob}_X(g\widetilde{x}_0) = \Phi_{\widetilde{X}}(g\widetilde{x})$，其中 $\Phi_{\widetilde{X}}$ 是绝对 Frobenius。由函子性，绝对 Frobenius 是带 $\pi_1^{\et}(X)$-作用的空间之间的映射，故 $\Phi_{\widetilde{X}}(g\widetilde{x}_0) = g \cdot \Phi_{\widetilde{X}}(\widetilde{x}_0)$。
  与我们先前的等式比较，可见该表达式等于 $F^{-1}\pi_1^{\et}(\Frob_X)(g) \cdot \widetilde{\Frob}_X(\widetilde{x}_0) = F^{-1}\pi_1^{\et}(\Frob_X)(g)F \cdot F^{-1}\widetilde{\Frob}_X(\widetilde{x}_0) = F^{-1}\pi_1^{\et}(\Frob_X)(g)F \cdot \Phi_{\widetilde{X}}(\widetilde{x}_0)$。因此 $g = \pi_1^{\et}(\Frob_X)(g)F$，断言得证。% [Tutor Note: 原文第二句作 “Φ_X̃(gx̃)”，据上下文应为 “Φ_X̃(gx̃_0)”]
  }，也就是说，我们有
\begin{equation}
  \pi_1^{\et}(\Frob_X)(g) = {}^F g.
\end{equation}

\pnum{点 $\widetilde{\Frob}_X(\widetilde{x}_0)$ 的选取还定义了一个同构
\[
  x_0^* \simeq \Frob_X(x_0)^* : \mathrm{Lisse}(X)^\heartsuit \to \Vect^\heartsuit
\]
（$\e$-线性对称幺半函子之间）。事实上，我们断言它来自对称幺半 DG 函子的同构
\begin{equation}
  x_0^* \simeq \Frob_X(x_0)^* : \qLisse(X) \to \Vect.
\end{equation}
确实，只要 $\qLisse(X)$ 是其核心的导出范畴，这就形式地成立。由 [2] 定理 E.2.8，$X \neq \mathbb{P}^1$ 时正是如此。稍初等一些（并且允许亏格 $0$），我们取 $U \subseteq X$ 为包含 $x_0$ 与 $\Frob_X(x_0)$ 的仿射开集，此时由（更简单的）定理 E.2.8 (a)，$\qLisse(U)$ 是其核心的导出范畴。}

\begin{remark}
下面的注不会被用到。眼下暂且让 $k$ 是\emph{任意}代数闭域（而不只是 $\overline{\F}_q$）。令 $Y/k$ 为有限型的连通概形，$y_1, y_2 \in Y(k)$ 为两个点。此时自然要问：是否存在对称幺半 DG 函子的同构
\[
  y_1^* \simeq y_2^* : \qLisse(Y) \to \Vect ?
\]
我们断言这是对的。确实，光滑连通曲线的情形我们上面已经有效地处理过了。任意连通曲线的情形可以按显然的方式归结为正规化（把奇异曲线各不可约分支之间的交点当作路标来引路）。最后，一般情形由如下初等论证得到：存在一条连通曲线 $C$ 与一个映射 $C \to Y$，使 $y_1$ 与 $y_2$ 都落在其像中。（人们不禁要问：在这个导出设定中是否存在一个纯 Tannaka 的论证？）
\end{remark}

\subsection{伴随商}

以下令 $H$ 为仿射代数群。

\pnum{令 $\LS^{\restr,\circ}_H$ 表示 $\LS^{\restr}_H$ 的\emph{单位}连通分支，即包含 $X$ 上平凡 $H$-局部系统的那个连通分支。我们提醒（见 [2] 命题 3.7.2）：$\LS^{\restr}_H$ 参数化（在某种确切的意义下）$X$ 上半单化平凡的 $H$-局部系统；下文中我们把这些称为\emph{幂幺} $H$-局部系统。

于是我们置
\[
  \LS^{\arthm,\circ}_H := \LS^{\arthm}_H \times_{\LS^{\restr}_H} \LS^{\restr,\circ}_H .
\]
换言之，$\LS^{\arthm,\circ}_H$ 是 Frobenius 作用在 $\LS^{\restr,\circ}_H$ 上的不动点；可以把它设想为参数化\emph{几何幂幺}的 Weil $H$-局部系统。}

\pnum{存在一个典范映射 $\mathbf{B}H \to \LS^{\restr,\circ}_H$，对应于平凡 $H$-局部系统。该映射是 Frobenius 等变的，其中 Frobenius 在 $\mathbf{B}H$ 上平凡地作用。

过渡到 Frobenius 不动点，我们得到映射
\begin{equation}
  \tau_H : H^{ad}/H \to \LS^{\arthm,\circ}_H .
\end{equation}}

\begin{remark}
粗略地说，复合 $H \to H^{ad}/H \to \LS^{\arthm,\circ}_H$ 把 $h \in H$ 送到由
\[
  \mathcal{W}_X \twoheadrightarrow \mathcal{W}_X/\pi_1^{\et}(X) \xrightarrow{\ \mathrm{Fr}\ } \ZZ \xrightarrow{\ 1 \mapsto h\ } H
\]
定义的 Weil 表示 $\mathcal{W}_X \to H$。由于我们只考虑 Weil 表示的\emph{群胚}，该映射如所愿地通过伴随商分解。
\end{remark}

\begin{remark}
下文的一个主要想法是：$\tau_H$ 离同构并不太远。为给后文提供动机，我们指出 $\tau_H$ 成为同构的如下障碍。粗略地说（例如在域值点的层面上），敌人显然是那些几何幂幺、且几何非平凡的 Weil 表示。

设 $H = GL_2$。令 $\lambda \in \e^\times$，并令 $\sigma_\lambda$ 表示 Frobenius 以乘 $\lambda$ 作用（而 $\pi_1^{\et}(X)$ 平凡作用）的一维 Weil 群表示。扩张 $0 \to \sigma_\lambda \to \sigma \to \sigma_1 \to 0$ 由 $\ZZ$ 的适当群上同调分类，即由复形
\[
  \underline{\Hom}_{\e[\ZZ]\text{-}\mathrm{mod}}\bigl(C_{\et,\bullet}(X), \sigma_\lambda\bigr)
\]
的 $H^1$ 分类。这里 $C_{\et,\bullet}(X)$ 是 $X$ 的\emph{\et{} 同调}的复形，$\ZZ$-作用的生成元在 \et{} 同调上按几何 Frobenius 作用。若 $\lambda \neq 1$，容易看出有正合列
\[
  0 \to H^1\bigl(\underline{\Hom}_{\e[\ZZ]\text{-}\mathrm{mod}}(C_{\et,\bullet}(X),\sigma_\lambda)\bigr) \to H^1_{\et}(X,\sigma_\lambda) \xrightarrow{\ \phi_X^{-1}-\id\ } H^1_{\et}(X,\sigma_\lambda),
\]
其中 $\phi_X$ 是作用在 $H^1_{\et}$ 上的几何 Frobenius。\footnote{注意，同调与上同调上的几何 Frobenius 是互为转置（即对偶）的态射。然而，如果我们把例如同调 $H_{\et,1}(X) = (\pi_1^{\et}(X)^{\mathrm{ab}})^{\wedge}_\ell \otimes \e$ 连同其几何 Frobenius（由 (4.1)，它对应于 $\Ad_F : \pi_1^{\et}(X) \to \pi_1^{\et}(X)$）看作 $\ZZ$-表示，那么上同调上的\emph{对偶} $\ZZ$-作用的生成元按\emph{算术} Frobenius 作用。这就解释了上式中出现的那个逆号。} 这里 $H^1_{\et}(X,\sigma_\lambda) = H^1_{\et}(X,\e)$，只是 Frobenius 作用取为标准作用的 $\lambda$ 倍。因此，若 $\lambda$ 是出现在 $H^1_{\et}(X,\e)$ 中的 Frobenius 特征值，我们就得到所需类型的\emph{几何}非平凡扩张。
\end{remark}

\pnumt{分裂。}{}通过限制到 $x_0$，我们得到映射 $\LS^{\restr}_H \to \mathbf{B}H$。由 §4.1.2\footnote{确切地说，当 $x_0$ 不是 $\F_q$-有理点时，这种等变性依赖于辅助选择。}，该映射把 Frobenius 与恒等映射缠结在一起，于是在不动点上我们得到映射
\[
  \LS^{\arthm}_H \to H^{ad}/H .
\]
我们把这个映射记作 $\chi_H = \chi_{H,x_0}$，并类似地记它在 $\LS^{\arthm,\circ}_H$ 上的限制。

由构造，复合
\[
  H^{ad}/H \xrightarrow{\ \tau_H\ } \LS^{\arthm}_H \xrightarrow{\ \chi_H\ } H^{ad}/H
\]
是恒等映射。
